\documentclass[11pt]{article}

\usepackage[margin=1in]{geometry}
\usepackage[T1]{fontenc}
\usepackage[utf8]{inputenc}
\usepackage{lmodern}
\usepackage{amsmath}
\usepackage{amssymb}
\usepackage{booktabs}
\usepackage{enumitem}
\usepackage{hyperref}

\hypersetup{
  colorlinks=true,
  linkcolor=blue,
  urlcolor=blue,
  pdftitle={Top Reconstruction in the Current Workspace},
  pdfauthor={Codex}
}

\title{Top Reconstruction in the Current Workspace}
\author{Prepared from Repository Material Only}
\date{\today}

\begin{document}

\maketitle

\begin{abstract}
This article summarizes the concept of top reconstruction, the concrete implementation present in this workspace, and the ways the resulting objects and scores are used downstream. The discussion is intentionally limited to material contained in the repository itself: the local \texttt{README} files, schema notes, run instructions, scripts, packaged pipeline code, and existing artifact reports. Within that scope, the workspace contains two closely related ideas. The first is a direct combinatorial reconstruction path in which jet triplets are enumerated and the highest-$p_T$ triplet is taken as a top candidate. The second is a modular machine-learning pipeline in which all candidate triplets are labeled using truth information, described by physics-motivated features, split into training and evaluation samples, scored by a classifier, and optionally converted back into event-level reconstructed top candidates.
\end{abstract}

\section{Introduction}

The central reconstruction problem described throughout this workspace is the identification of a hadronic top candidate from jets. The slide deck \texttt{top\_reco\_slides.tex} states the physics picture as
\[
t \to Wb \to qq'b,
\]
and emphasizes the combinatorial difficulty that arises once an event contains more than three jets. If an event has $N$ jets, then there are $\binom{N}{3}$ unordered three-jet combinations. Only a subset of those combinations, and often only one of them, corresponds to the desired top decay pattern. The repository therefore frames top reconstruction as a structured selection problem over triplets.

The workspace also makes clear that this problem is studied in a practical software setting rather than only as a conceptual exercise. The code and documentation define a pipeline that starts from ROOT TTrees, computes triplet-level kinematic features, writes Parquet datasets, trains either an XGBoost classifier or a TabPFN model, runs inference on all candidate triplets, and then uses those scores to build event-level top-candidate outputs. Alongside that pipeline, several smaller scripts support exploratory plotting, cutflow studies, and validation of branch semantics.

\section{Conceptual View of Top Reconstruction}

At the conceptual level, the workspace supports two complementary reconstruction viewpoints.

The first viewpoint is purely combinatorial. The script \texttt{triplet\_reco.py} takes a per-event jet array and evaluates every unordered triplet of non-padded jets. For each combination it forms the summed four-vector, computes the triplet transverse momentum, pseudorapidity, azimuth, and invariant mass, and chooses the single triplet with the largest system $p_T$. In this formulation, top reconstruction is an optimization problem over candidate subsets. The output is not a class probability but a chosen triplet and the distributions of its kinematic observables.

The second viewpoint is statistical and supervised. The schema note \texttt{schema.md} defines a triplet-level dataset in which every unordered triplet is treated as an example and is labeled as either truth-matched or not truth-matched. The classifier then learns to distinguish ``signal'' triplets, meaning triplets that match a truth definition, from ``background'' triplets, meaning all other combinations. In this formulation, top reconstruction is a binary classification problem defined on candidate triplets. Event-level reconstruction can then be recovered by ranking and selecting the scored triplets.

These two viewpoints are not contradictory. Rather, the simpler highest-$p_T$ triplet path gives an intuitive baseline, while the machine-learning path attempts to replace a single heuristic ranking variable with a richer score informed by internal triplet geometry and mass structure.

\section{Input Representation and Truth Information}

The schema file and the \texttt{dataset\_build.py} implementation define the expected ROOT inputs. The required branches are
\begin{itemize}[leftmargin=*]
  \item \texttt{N\_genjet},
  \item \texttt{genjet\_pt}, \texttt{genjet\_eta}, \texttt{genjet\_phi},
  \item \texttt{truth\_triplet\_0} through \texttt{truth\_triplet\_3},
  \item and an event identifier branch, with \texttt{Number} used by default when it exists.
\end{itemize}

The schema explicitly states that jets are treated as massless for the feature-building stage and that any genjet mass branch must be ignored. The event identifier is expected to come from the branch interpretation documented elsewhere in the repository, but the code also contains a fallback to the entry index if the requested event-number branch is absent.

Truth information is represented as up to four triplets per event, stored in the \texttt{truth\_triplet\_*} branches. The implementation parses each branch as a length-three integer array, rejects malformed entries, rejects repeated indices, and rejects indices outside the range $0 \le i < N_{\mathrm{genjet}}$. A key design choice is that truth triplets are treated as unordered sets. In practice this means that the triplet $(1,4,7)$ is equivalent to $(7,1,4)$ after sorting. The build stage records both the total number of valid truth triplets and the number of examples whose original branch ordering was not already sorted.

\section{Candidate Construction and Physics Features}

The candidate-building rule in \texttt{dataset\_build.py} is deliberately exhaustive. For every event with $N_{\mathrm{genjet}}$ jets, the code loops over all unordered combinations
\[
i < j < k
\]
and constructs a candidate row for every such triplet. The code checks that the produced number of triplets matches the combinatorial expectation $\binom{N_{\mathrm{genjet}}}{3}$.

Each candidate is described by six engineered features, defined in \texttt{features.py} and summarized in \texttt{schema.md}. Three of them are angular distances:
\begin{align}
\Delta R_{ab} &= \sqrt{(\eta_a - \eta_b)^2 + \Delta \phi_{ab}^2}, \\
\Delta R_{ac} &= \sqrt{(\eta_a - \eta_c)^2 + \Delta \phi_{ac}^2}, \\
\Delta R_{bc} &= \sqrt{(\eta_b - \eta_c)^2 + \Delta \phi_{bc}^2},
\end{align}
where the azimuthal difference is wrapped into the principal range through the usual trigonometric construction.

The other three features are pair-mass fractions:
\begin{align}
\frac{m_{ab}}{m_{123}}, \qquad
\frac{m_{ac}}{m_{123}}, \qquad
\frac{m_{bc}}{m_{123}}.
\end{align}
To compute these quantities, the code builds massless jet four-vectors according to
\begin{align}
p_x &= p_T \cos \phi, &
p_y &= p_T \sin \phi, &
p_z &= p_T \sinh \eta, &
E &= \sqrt{p_x^2 + p_y^2 + p_z^2},
\end{align}
and then uses
\[
m^2 = E^2 - |\vec{p}|^2.
\]
Any negative value of $m^2$ caused by numerical precision is clamped to zero before taking the square root. If the full triplet mass $m_{123}$ is zero, the ratio features are set to zero.

In addition to the six training features, the code records several observables that are useful for downstream diagnostics: the triplet invariant mass $m_{123}$, the three pairwise masses $m_{ab}$, $m_{ac}$, and $m_{bc}$, and the reconstructed triplet kinematics $(p_T,\eta,\phi)$. The label column \texttt{is\_truth} is set to one when the sorted candidate indices belong to the event's truth set and zero otherwise.

\section{A Direct Reconstruction Path: Cutflow and Highest-$p_T$ Triplets}

Before introducing the machine-learning pipeline, the repository offers a more direct reconstruction path built around \texttt{cutflow\_and\_store.py} and \texttt{triplet\_reco.py}. This path is useful because it makes the reconstruction logic visible with minimal infrastructure.

According to the root-level \texttt{README.md}, \texttt{cutflow\_and\_store.py} reads a ROOT file, applies a simple jet selection, and stores selected event content for events with at least six selected jets. The default jet requirements are
\[
p_T > 25~\mathrm{GeV}, \qquad |\eta| < 2.5.
\]
The script writes a cutflow table and a padded event array \texttt{selected\_jets.npy} with shape $(N_{\mathrm{selected\ events}}, 10, 4)$, where the last axis stores $(p_T,\eta,\phi,m)$. It also writes a richer archive \texttt{selected\_events.npz} that includes truth triplet indices, leading-triplet kinematics, up to four disjoint top-candidate slots, and a placeholder \texttt{top\_PID} array.

The \texttt{README} further explains that the stored triplet extraction is sequential and disjoint: the code finds the highest-$p_T$ three-jet system, removes those jets, and repeats the procedure up to four times. This is already a useful event-level reconstruction heuristic because it turns a variable number of jets into a small, structured set of candidate top objects.

The companion script \texttt{triplet\_reco.py} then reads \texttt{selected\_jets.npy}, ignores padded jets with zero transverse momentum, evaluates all unordered three-jet combinations for each event, and keeps the candidate with the largest system $p_T$. It produces histograms of the selected triplet's $p_T$, $\eta$, $\phi$, and mass, plus a simple HTML index page. In other words, this direct reconstruction path turns selected jets into a single top-like candidate per event using a transparent ranking rule.

\section{Validation of Branch Interpretation}

The repository also contains \texttt{validate\_triplet\_interpretation.py}, which addresses an important practical problem: the semantic meaning of triplet-related branches in the input ntuple. That script compares several alternative interpretations, including applying \texttt{reco\_triplet\_0} or \texttt{truth\_triplet\_0} to reconstructed jets or generator-level jets, and then overlaying the resulting kinematic distributions against truth-level top branches \texttt{top\_*}. The presence of this validation tool shows that reconstruction in this workspace is not only about training a classifier; it also requires careful checking that branch names and index conventions are understood correctly before the outputs are trusted.

\section{The Modular Machine-Learning Pipeline}

The main implementation described in \texttt{schema.md}, \texttt{run\_intruction.md}, \texttt{clean\_repo/README.md}, and the Python modules follows a staged workflow. Before naming the stages in software terms, it is useful to describe them conceptually.

\begin{itemize}[leftmargin=*]
  \item \textbf{Stage 1 --- generate the physically allowed top hypotheses.} At the conceptual level, each event is treated as a collection of possible three-jet top-decay interpretations. The goal is not to choose immediately, but to enumerate the candidate hypotheses and mark which ones agree with the truth definition. In the codebase, this is the \texttt{dataset\_build} stage: ROOT input is converted into \texttt{triplets\_raw.parquet}, each row carries \texttt{event\_id}, triplet indices \texttt{i,j,k}, the engineered features, observables, and the label \texttt{is\_truth}.
  \item \textbf{Stage 2 --- organize the learning problem without breaking event coherence.} Conceptually, one wants to learn from some events, validate on others, and test on a separate sample, while keeping all triplets from the same event together. At the same time, one wants to prevent the many fake triplets from overwhelming the rare truth-matched ones. In the codebase, this is \texttt{dataset\_prepare}: the stage applies deterministic event-level splitting, uses \texttt{assign\_split}, applies the \texttt{max\_bg\_per\_event} cap to train and validation, and writes \texttt{train.parquet}, \texttt{val.parquet}, and \texttt{test.parquet}.
  \item \textbf{Stage 3 --- learn what a real top-like triplet looks like.} In physics language, the model is learning the internal kinematic pattern of a genuine top candidate as distinct from accidental jet combinations. In machine-learning language, this is supervised binary classification of truth versus combinatorial background. In the codebase, this is the \texttt{train} stage in which either \texttt{xgb} or \texttt{tabpfn} is trained from the six \texttt{FEATURE\_COLUMNS}, and the outputs include \texttt{model\_xgb.json} or \texttt{model\_tabpfn.pkl} together with \texttt{training\_report*.json}.
  \item \textbf{Stage 4 --- assign a plausibility score to every candidate.} Conceptually, this stage turns the trained model into a ranking function over all remaining top hypotheses. Rather than selecting only one candidate immediately, it preserves the full candidate population together with a learned measure of how signal-like each triplet is. In the codebase, this is \texttt{infer}: the stage reads the test dataset in batches, appends a score column such as \texttt{score} or \texttt{score\_tabpfn}, and writes \texttt{inference\_test.parquet}.
  \item \textbf{Stage 5 --- turn triplet scores into event-level reconstructed tops.} Conceptually, the reconstruction problem returns to the event level: the scored triplets are used to select one or more top candidates that summarize the event in a physically interpretable way. In the codebase, this is the packaged \texttt{select\_triplets} stage under \texttt{clean\_repo/src/triplet\_ml}: strategies such as \texttt{greedy\_disjoint}, \texttt{top1}, and \texttt{best\_pair\_avg\_disjoint} write \texttt{selected\_triplets.parquet} and \texttt{event\_selection.parquet}.
\end{itemize}

This staged design matters because it decouples expensive combinatorial enumeration from model development. Once \texttt{triplets\_raw.parquet} has been produced, different balancing settings, different models, and repeated training runs can be attempted without re-reading the ROOT input and without re-enumerating all triplets.

\subsection{Stage 1: Building the Triplet Dataset}

Conceptually, the first stage turns a single event into a full space of top-candidate hypotheses. It does this by refusing to guess too early: every allowed three-jet interpretation is kept, and the truth information is used only to label which hypotheses are physically correct.

Technically, the \texttt{dataset\_build.py} stage uses \texttt{uproot.iterate} to stream over the required ROOT branches. It resolves the TTree name automatically unless the user provides one explicitly. For each event it determines the event identifier, copies the generator-level jet arrays, constructs the event's truth-triplet set, and loops over all candidate triplets. Every candidate row is appended into a column buffer and periodically flushed to a streaming Parquet writer. This is why the schema describes \texttt{dataset\_build} as the expensive stage.

The outputs of this stage are more than just the Parquet file. A JSON report records, among other things, the number of processed events, the number of candidate triplets written, the total number of valid truth triplets, and the event-id branch that was actually used. A separate \texttt{config\_snapshot.json} captures the input files, tree name, requested event-id branch, maximum number of events, chunk size, and row-group settings.

\subsection{Stage 2: Event-Level Splitting and Balancing}

Conceptually, the second stage prepares a fair learning problem. It preserves the idea that an event is a coherent object while also preventing the overwhelming population of fake triplets from dominating the statistical picture.

Technically, the \texttt{dataset\_prepare.py} stage reads the raw triplet Parquet dataset, groups rows by \texttt{event\_id}, and assigns each event to train, validation, or test using a deterministic hash-based split. The default fractions are
\[
20\% \text{ train}, \qquad 10\% \text{ validation}, \qquad 70\% \text{ test}.
\]
Because the split is assigned at the event level, all triplets from the same event remain in the same partition, preventing leakage across datasets.

This stage also implements a hybrid treatment of class imbalance. The repository documentation notes that true triplets are rare compared with combinatorial backgrounds. To manage this, the code caps the number of background triplets per event in the train and validation splits while leaving the test split untouched. The default background cap is 200 triplets per event. The kept background rows are chosen deterministically using a stable hash of the event and candidate indices, so rerunning with the same seed preserves the same retained subset.

In addition to writing \texttt{train.parquet}, \texttt{val.parquet}, and \texttt{test.parquet}, the stage writes both machine-readable and human-readable summaries of the split composition. The resulting reports include the number of events per split, the number of truth and background triplets, the number removed by balancing, and average triplets per event.

\subsection{Stage 3: Model Training}

Conceptually, the third stage asks the model to discover which internal triplet patterns are characteristic of real top decays. The classifier is therefore not learning event identity; it is learning a discriminant over candidate triplets.

Technically, the training stage, implemented in \texttt{train.py}, supports two backends: XGBoost and TabPFN. In both cases, the training code loads exactly the six engineered feature columns together with the truth label. The current code computes internal summary metrics such as binary AUC and validation log loss, saves the trained model, and writes a training report and a short Markdown statistics summary.

For XGBoost, the repository exposes explicit hyperparameters through the command-line interface. The defaults in the current code are
\begin{itemize}[leftmargin=*]
  \item \texttt{num\_boost\_round} = 400,
  \item \texttt{early\_stopping\_rounds} = 30,
  \item learning rate \texttt{eta} = 0.05,
  \item \texttt{max\_depth} = 6,
  \item \texttt{min\_child\_weight} = 1.0,
  \item \texttt{subsample} = 0.8,
  \item \texttt{colsample\_bytree} = 0.8,
  \item and \texttt{tree\_method} = \texttt{hist}.
\end{itemize}
An optional \texttt{--use-sample-weights} switch applies class-imbalance weighting by upweighting the positive class according to the train and validation class counts.

For TabPFN, the code supports a separate path in which the training sample can be explicitly subsampled and class-balanced to a $50/50$ signal-background mixture. This is controlled by \texttt{--max-training-samples} and \texttt{--tabpfn-balance-classes}. The resulting report captures not only metrics but also the sampling strategy that was used to form the final training subset.

\subsection{Stage 4: Inference on All Candidate Triplets}

Conceptually, the fourth stage turns learned knowledge into a ranking over all candidate top interpretations. Instead of collapsing the event immediately to one answer, it first scores the entire candidate population.

Technically, the \texttt{infer.py} stage loads a trained model, reads the test Parquet file in batches, reconstructs the feature matrix, and appends a score column to each row before writing the scored dataset back to Parquet. The row count is checked at the end to ensure that the number of input triplets matches the number of output scores.

This design has an important consequence for later analysis. The stage does not discard low-score triplets or attempt immediate event-level reconstruction. Instead, it preserves the full candidate set. This makes the scored Parquet file a reusable intermediate product: different score thresholds, different event-level selection strategies, and different plotting studies can all be applied downstream without rerunning inference.

The plotting code connected to the inference stage generates ROC curves, score distributions, threshold scans, and, when applicable, score-versus-mass views. Comparison plots between XGBoost and TabPFN are also supported when both scored outputs are present.

\subsection{Stage 5: Turning Scores into Event-Level Top Candidates}

Conceptually, the fifth stage closes the loop back to the physics object one actually wants: a small number of top candidates per event. At this point the score is no longer just a diagnostic quantity; it becomes the basis for reconstruction decisions.

Technically, the packaged repository under \texttt{clean\_repo/src/triplet\_ml} adds a fifth stage, \texttt{select\_triplets.py}, which makes the role of the classifier especially clear. At this stage, the triplet score is no longer an end in itself; it becomes a ranking variable used to build event-level top-candidate outputs.

The code supports multiple strategies:
\begin{itemize}[leftmargin=*]
  \item \texttt{greedy\_disjoint},
  \item \texttt{top1},
  \item \texttt{topk},
  \item \texttt{threshold},
  \item and \texttt{best\_pair\_avg\_disjoint}.
\end{itemize}

The stage writes two Parquet outputs. The first, \texttt{selected\_triplets.parquet}, stores the selected candidates themselves together with their rank and score. The second, \texttt{event\_selection.parquet}, stores event-level summaries such as the number of selected top candidates, fixed top-candidate slots \texttt{top1\_*} through \texttt{top4\_*}, and the invariant mass of the two leading selected candidates. In the pair-based strategy, the code explicitly requires enough inferred jets to define two disjoint candidates and leaves dummy values when no acceptable pair is available.

This stage explains how the machine-learning outputs are used in practice. The classifier does not itself reconstruct a final top object; it produces a scoring function. Event-level top reconstruction is then recovered by applying a strategy that translates scored triplets into a small number of candidate tops per event.

\section{Interfaces and Workflow in Practice}

The repository contains two parallel interfaces to essentially the same logic. At the root level, a script-style CLI built around \texttt{main.py} can run the individual stages directly. In parallel, the packaged layout in \texttt{clean\_repo} exposes the module \texttt{triplet\_ml} and a console entry point, and it also provides a \texttt{Makefile} for common end-to-end workflows.

The run instructions found in the repository consistently describe the core data flow as
\[
\texttt{ROOT} \rightarrow \texttt{triplets\_raw.parquet} \rightarrow
\{\texttt{train,val,test}\} \rightarrow \texttt{model} \rightarrow \texttt{inference parquet}.
\]
For exploratory work, the user can instead begin with \texttt{plot\_jets.py}, \texttt{cutflow\_and\_store.py}, and \texttt{triplet\_reco.py}. For model-building work, the user typically proceeds through the staged Parquet-based pipeline. For scored event reconstruction studies, the user continues from inference into \texttt{select\_triplets.py}.

\section{Examples from Existing Workspace Artifacts}

The repository already contains several example run directories, which help illustrate the scale of the data products produced by the pipeline. Table~\ref{tab:workspace-runs} summarizes a few of them using the JSON reports currently stored in the workspace.

\begin{table}[htbp]
  \centering
  \caption{Example artifact summaries already present in the workspace.}
  \label{tab:workspace-runs}
  \begin{tabular}{@{}llll@{}}
    \toprule
    Run directory & Input sample & Processed events & Triplet rows \\
    \midrule
    \texttt{artifacts/run\_10000\_fresh} & \texttt{ttbar.root} & 10{,}000 & 105{,}784 \\
    \texttt{artifacts/run\_40000} & \texttt{ttbar.root} & 40{,}000 & 419{,}728 \\
    \texttt{study/artifacts/run\_last10k} & \texttt{study/data/ttbar\_last\_10000.root} & 10{,}000 & 104{,}892 \\
    \bottomrule
  \end{tabular}
\end{table}

The split reports present in the workspace also show the effect of deterministic event splitting and class imbalance. For example, the \texttt{run\_10000\_fresh} prepare stage records 21{,}630 training rows, 11{,}953 validation rows, and 72{,}201 test rows, while the corresponding inference output contains 72{,}201 scored triplets. The \texttt{run\_40000} directory shows the same pattern at larger scale, with 292{,}017 triplets in the test set and the same number of scored rows after inference. These reports demonstrate that the workspace already uses the staged design in concrete runs rather than only in abstract documentation.

\section{What the Reconstruction Outputs Mean}

Across the workspace, reconstruction outputs appear in several forms, each serving a different purpose.

First, there are \emph{event-level jet snapshots} such as \texttt{selected\_jets.npy} and \texttt{selected\_events.npz}. These are useful for direct combinatorial studies and for visual inspection of selected events.

Second, there are \emph{triplet-level learning datasets} such as \texttt{triplets\_raw.parquet}, \texttt{train.parquet}, \texttt{val.parquet}, and \texttt{test.parquet}. These are the central objects used for supervised learning. They preserve candidate indices, features, observables, and truth labels.

Third, there are \emph{scored triplet datasets} such as \texttt{inference\_test.parquet}. These are crucial because they retain the full candidate set while appending a learned ranking variable.

Finally, there are \emph{event-level reconstructed top selections} such as \texttt{selected\_triplets.parquet} and \texttt{event\_selection.parquet}. These files represent the point at which triplet scores are converted back into a physically interpretable collection of candidate top objects per event.

In short, the workspace uses top reconstruction outputs in two ways. They are used diagnostically, through histograms, validation overlays, and cutflow summaries, and they are used operationally, as inputs to later ranking, selection, and event-structure studies.

\section{Limitations and Design Choices Visible in the Workspace}

Several design choices are explicit in the local code and documentation.

The machine-learning build stage uses all unordered generator-level jet triplets and does not apply explicit $p_T$, $\eta$, or $b$-tag requirements at candidate-construction time. This keeps the dataset definition simple and exhaustive, but it also means that the classifier must absorb a large combinatorial background.

The feature builder treats jets as massless and uses only six engineered variables for training. This is physically motivated and compact, but it intentionally limits the feature space compared with a more elaborate object-level representation.

The deterministic split-and-cap logic is designed for reproducibility, yet the repository notes one practical caveat: very small event samples can produce empty validation sets and therefore fail training. The workspace also contains historical artifact directories with slightly different Parquet schemas, which is a reminder that pipeline outputs evolved over time.

These limitations do not undermine the utility of the repository. Rather, they define its scope. The current workspace is best understood as a well-structured study environment for triplet-based top reconstruction, with both transparent baseline methods and a more reusable ML-centered workflow.

\section{Conclusion}

Based on the files present in this workspace, top reconstruction is treated as the problem of selecting physically meaningful three-jet combinations from a larger combinatorial set. The repository supports a direct highest-$p_T$ triplet baseline, validation tools for understanding branch semantics, and a modular machine-learning pipeline that converts ROOT TTrees into scored triplet datasets and then, if desired, back into event-level top candidates.

What makes the implementation especially coherent is the separation between stages. Candidate generation, dataset preparation, model training, inference, and event-level selection are all isolated into reusable steps with explicit artifacts and reports. This allows the same underlying reconstruction question to be studied from several angles: as a cutflow problem, as a combinatorial ranking problem, as a supervised learning task, and as an event-level object-building exercise.

\appendix

\section{Workspace Basis and Scope}

This note was written only from material found in the current workspace. The key local sources used were:
\begin{itemize}[leftmargin=*]
  \item \texttt{top\_reco\_slides.tex},
  \item \texttt{README.md},
  \item \texttt{REPO\_SUMMARY.md},
  \item \texttt{schema.md},
  \item \texttt{run\_intruction.md},
  \item \texttt{features.py},
  \item \texttt{dataset\_build.py},
  \item \texttt{dataset\_prepare.py},
  \item \texttt{train.py},
  \item \texttt{infer.py},
  \item \texttt{triplet\_reco.py},
  \item \texttt{validate\_triplet\_interpretation.py},
  \item \texttt{clean\_repo/README.md},
  \item \texttt{clean\_repo/src/triplet\_ml/select\_triplets.py},
  \item and the JSON reports already stored under \texttt{artifacts/} and \texttt{study/artifacts/}.
\end{itemize}

No external physics references, detector notes, or online documentation were used in preparing this article.

\end{document}
